% %%%

%%%   Самарский государственный технический университет

%%%   Кафедра прикладной математики и информатики

%%%

%%%   Преамбула для статьи в журнал "Вестник Самарского государственного

%%%   технического университета. Серия: «Физико-математические науки»"

%%%   Ориентирована под MikTEx 2.4 for Win32

%%%

%%%   Хотя сам журнал собирается под GNU/Linux на дистрибутиве TeXLive



\documentclass[11pt,a4paper,twoside]{article}



\usepackage[cp1251]{inputenc}

\usepackage[T2A]{fontenc}

\usepackage[english,russian]{babel}

\usepackage{samgtubib}

\usepackage[dvips]{graphicx}

\usepackage[multiple]{footmisc}



\usepackage{samgtu}

\renewcommand{\VestnikYear}{2009} % Год издания Вестника

\renewcommand{\VestnikNumber}{2\,(19)} % Номер Вестника

\renewcommand{\VestnikMonth}{Ноябрь} % Месяц выхода Вестника

\renewcommand{\VestnikData}{01.11.09} % Дата подписания Вестника в печать

\renewcommand{\VestnikUPL}{26,64} % Усл. печ. л.

\renewcommand{\VestnikUIL}{25,73} % Уч.-изд. л.

\renewcommand{\VestnikRegNumber}{479/09} % Регистрационный номер

\renewcommand{\VestnikOrder}{994} % Номер заказа







%

%\usepackage{pst-barcode}

%\usepackage{auto-pst-pdf}

%%

%%\newcommand{\totop}[1]{\raisebox{6pt}[1pt][2pt]{#1}}

%%\newcommand{\tottop}[1]{\raisebox{12pt}[1pt][2pt]{#1}}

%%\newcommand{\totttop}[1]{\raisebox{18pt}[1pt][2pt]{#1}}

%%\newcommand{\tottttop}[1]{\raisebox{24pt}[1pt][2pt]{#1}}

%%\newcommand{\totttttop}[1]{\raisebox{30pt}[1pt][2pt]{#1}}

%%\newcommand{\tobot}[1]{\raisebox{-6pt}[1pt][2pt]{#1}}

%%\newcommand{\tobbot}[1]{\raisebox{-12pt}[1pt][2pt]{#1}}

%%\newcommand{\tobbbot}[1]{\raisebox{-18pt}[1pt][2pt]{#1}}

%

%\graphicspath{{./00PIC/}}

%

%%\samgtupubl % для генерации pdf, отдаваемого в типографию СамГТУ

%\pdfcrop % обрезка pdf для размещения на math-net.ru

%\draft % На корректуре

%\filllastpage % Последняя страница не заполнена не полностью

%%\briefcomm % Статья является кратким сообщением

%



\vsgtu{1386}



\FirstPage{44}

\LastPage{62}

%

%

%\begin{document}



\hfill \begin{pspicture}(1in, 1in)

\psbarcode{http://dx.doi.org/10.14498/vsgtu1386}{}{qrcode}

\end{pspicture}

\vspace{-1in}





\udk{517.956.25}

\msc{35J62, 35J25, 35J15}



\rutitle{О решениях эллиптических уравнений с~нестепенными нелинейностями в~неограниченных~областях}

{О решениях эллиптических уравнений с нестепенными нелинейностями \dots}



\titlehack{О решениях эллиптических уравнений с~нестепенными нелинейностями в~неограниченных~областях$^*$}



\ruauthor{Л.~М.~Кожевникова, А.~А.~Хаджи}

 {К\,о\,ж\,е\,в\,н\,и\,к\,о\,в\,а~Л.~М.,\ Х\,а\,д\,ж\,и~А.~А.}



\ruaffil{\bashgusterl.}



%\email{1}{\mem{kosul@gmail.ru}%, anna_5955@mail.ru}

%(L.M.~Kozhevnikova, \CAut)}





\ruauthordetails{2}{\fullauthor{\rupid{8695}{Лариса Михайловна Кожевникова}}\regalia{д.ф.-м.н., доц.; \mem{kosul@gmail.ru}; автор, ведущий переписку} \position{профессор} \dept{каф. математического анализа}\\

\fullauthor{\rupid{76306}{Анна Александровна Хаджи}} \regalia{\mem{anna_5955@mail.ru}} \position{старший преподаватель} \dept{каф. общенаучных  дисциплин}

\par

\smallskip

\noindent

$^*$Настоящая статья представляет собой расширенный вариант доклада \cite{kozh:add1-ru}, сделанного автором на Четвёртой международной конференции «Математическая физика и её приложения» (Россия, Самара, 25 августа -- 1 сентября 2014).

}



\enauthordetails{2}{\fullauthor{\enpid{8695}{Larisa M. Kozhevnikova}}\regalia{Dr. Phys. \& Math. Sci.; \mem{kosul@gmail.ru}; Corresponding Author}\position{Professor} \dept{Dept. of Mathematical Analysis}\\

\fullauthor{\enpid{76306}{Anna A. Khadzhi}} \regalia{\mem{anna_5955@mail.ru}}  \position{Senior Teacher} \dept{Dept. of Scientific Disciplines}

\par

\smallskip

\noindent

$^*$This paper is an extended version of the paper \cite{kozh:add1-eng}, presented at the Mathematical Physics and Its Applications 2014 Conference.}



\rukeywords{анизотропное эллиптическое уравнение, пространство Соболева"--~Орлича,

нестепенные нелинейности, существование решения, неограниченная область, ограниченность решения, убывание решения}



\ruabstract{В работе выделен некоторый класс анизотропных эллиптических уравнений второго порядка дивергентного вида с младшими членами

с нестепенными нелинейностями

 $$\sum\limits_{\alpha=1}^{n}(a_{\alpha}({\bf x},u,\nabla u))_{x_{\alpha}}-a_0({\bf x},u,\nabla u)=0.$$

 На каратеодориевы функции, входящие в уравнение, накладывается  ус\-ловие  совокупной монотонности. Ограничения на  рост функций формулируются в терминах специального класса выпуклых функций. Эти требования обеспечивают ограниченность, коэрцитивность, монотонность и~семинепрерывность соответствующего эллиптического оператора. Для рассматриваемых уравнений с нестепенными нелинейностями исследованы качественные свойства решений задачи Дирихле в неограниченных областях $\Omega\subset \mathbb{R}_n,\;n\geq 2$. Установлены существование и единственность обобщённых решений    в анизотропных пространствах Соболева"--~Орлича.

 Кроме того,  для произвольных неограниченных областей обобщены теоремы вложения анизотропных пространств Соболева"--~Орлича. Это позволило доказать глобальную ограниченность  решений задачи Дирихле.  Использована оригинальная геометрическая характеристика для неограниченных областей, расположенных вдоль выделенной оси. В терминах этой характеристики установлена экспоненциальная оценка скорости убывания на бесконечности решений рассматриваемой задачи с финитными данными.}



\entitle{On solutions of elliptic equations with nonpower nonlinearities in unbounded domains}



\engtitlehack{On solutions of elliptic equations with nonpower nonlinearities in unbounded domains$^*$}



\enauthor{L.~M.~Kozhevnikova, A.~A.~Khadzhi}{K\,o\,z\,h\,e\,v\,n\,i\,k\,o\,v\,a~L.~M.,

K\,h\,a\,d\,z\,h\,i~A.~A.}



\enaffil{\engbashgusterl.}





\enabstract{The paper highlighted some class of anisotropic elliptic equations of second order in divergence form with younger members with nonpower nonlinearities 

$$\sum\limits_{\alpha=1}^{n}(a_{\alpha}({\bf x},u,\nabla u))_{x_{\alpha}}-a_0({\bf x},u,\nabla u)=0.$$

The condition of total monotony is imposed on the Caratheodory functions included in the equation. Restrictions on the growth of the functions are formulated in terms of a special class of convex functions. These requirements provide limited, coercive, monotone and semicontinuous corresponding elliptic operator. For the considered equations with nonpower nonlinearities the qualitative properties of solutions of the Dirichlet problem in unbounded domains $ \Omega \subset \mathbb {R} _n, \; n \geq 2$ are studied. The existence and uniqueness of generalized solutions in anisotropic Sobolev--Orlicz spaces are proved. Moreover, for arbitrary unbounded domains, the Embedding theorems for anisotropic Sobolev--Orlicz spaces are generalized. It makes possible to prove the global boundedness of solutions of the Dirichlet problem. The original geometric characteristic for unbounded domains along the selected axis is used. In terms of the characteristic the exponential estimate for the rate of decrease at infinity of solutions of the problem with finite data is set.}



\enkeywords{anisotropic elliptic equations, Sobolev--Orlicz space, nonpower nonlinearity, the existence of solution, unbounded domains, boundedness of solutions, decay of solution}



\date{15/XII/2014} % Дата поступления статьи в редакцию.

\revisiondate{13/II/2015} % В окончательном виде

\accepteddate{25/II/2015}



\makerutitle



\smallskip



\Section[n]{Введение}

Пусть $\Omega $ "--- произвольная неограниченная область пространства

 $\mathbb{R}_{n}=\{{\bf x}=(x_{1},x_{2},\dots,x_{n})\}$, $\Omega\subset\mathbb{R}_{n}$, $n \geq 2$.

 Для анизотропного квазилинейного эллиптического уравнения второго порядка рассматривается задача Дирихле

\begin{gather}\label{eq:Kozh:ur}

 \sum\limits_{\alpha=1}^{n}(a_{\alpha}({\bf x},u,\nabla u))_{x_{\alpha}}-a_0({\bf x},u,\nabla u)=0,\quad {\bf x} \in \Omega;

\\

\label{eq:Kozh:gu}

  u\big|_{\partial\Omega}= 0.

\end{gather}

Предполагается, что функции

$a_{\alpha}({\bf x},p_0,p)$, $\alpha=0,\ldots,n$,   измеримы по ${\bf x}\in \Omega$ для ${\bf p}=(p_0,p)=(p_0,p_1,\ldots,p_{n})\in\mathbb{R}_{n+1}$, непрерывны по ${\bf p}\in\mathbb{R}_{n+1}$ для почти всех ${\bf x}\in\Omega$. Ограничения на нелинейные функции $a_{\alpha}({\bf x},p_0,p)$, $\alpha=0,\ldots,n$, формулируются в терминах $N$-функций, они   будут приведены ниже.





В работе исследуются вопросы существования, ограниченности и убывания на бесконечности решений задачи \eqref{eq:Kozh:ur}, \eqref{eq:Kozh:gu}  в неограниченных областях~$\Omega$.



Существование решений для уравнений вида \eqref{eq:Kozh:ur} в ограниченных областях изучалось в работах [\citen{KozhKhad:bib:Vichik}--\citen{KozhKhad:bib:Klimov2}]. В нашей работе на  каратеодориевы функции, входящие в уравнение, наряду с условием монотонности наложены требования, которые позволяют установить существование единственного обобщённого решения задачи \eqref{eq:Kozh:ur}, \eqref{eq:Kozh:gu} в произвольной неограниченной области~$\Omega\subset\mathbb{R}_{n}$.



После нового доказательства теоремы Э.~де~Джорджи \cite{KozhKhad:bib:Giorgi} при помощи итерационной техники, данного Ю.~Мозером в работе \cite{KozhKhad:bib:Moser},   вопросы ограниченности и непрерывности по Гёльдеру обобщённых решений различных классов линейных и квазилинейных эллиптических уравнений с изотропными степенными нелинейностями исследовались в работах С.\,Н.~Кружкова \cite{KozhKhad:bib:Kruzhkov}, Дж.~Серрина \cite{KozhKhad:bib:Serrin}, Е.~М.~Ландиса \cite{KozhKhad:bib:Landis} и других авторов.





И.~М.~Колодий \cite{KozhKhad:bib:Kolodii} установил ограниченность решений некоторого класса анизотропных эллиптических уравнений со степенными нелинейностями в ограниченных областях. При этом требование ограниченности области является существенным условием в его доказательстве.  Для неограниченных областей этот результат был получен в работе \cite{KozhKhad:bib:KozhKhad2}. Попытка А.~Г.~Королёва \cite{KozhKhad:bib:Korolev3} осуществить дальнейшее распространение техники Мозера на эллиптические дифференциальные уравнения с анизотропными  нестепеными нелинейностями  содержит досадную ошибку.



Развивая метод априорных оценок для срезок \cite[гл. II, \S~5, лемма 5.1]{KozhKhad:bib:LadizhYralc},

В.~С.~Климов в работе \cite{KozhKhad:bib:Klimov3} для некоторого вида изотропных эллиптических уравнений с нестепенными нелинейностями доказал глобальную ограниченность решений в ограниченных областях.

А.~Г.~Королёв  \cite{KozhKhad:bib:Korolev2} получил интегральный вариант теоремы вложения для функций из пространства Соболева"--~Орлича, на его основе

он доказал ограниченность решений для некоторого класса анизотропных эллиптических  уравнений с нестепенными нелинейностями в ограниченных областях. Авторы  настоящей статьи  наложили требования на структуру уравнения, позволившие  установить ограниченность решений в неограниченных  областях.



Изучению поведения на бесконечности решений линейных

эллиптических уравнений в неограниченных областях посвящены работы О.~А.~Олейник, Г.~А.~Иосифьяна \cite{KozhKhad:bib:OlIosif}, Е.~М.~Ландиса, Г.~П.~Панасенко, В.~А.~Кондратьева, И.~Копачека, Д.~М.~Леквеишвили, О.~А.~Олейник \cite {KozhKhad:bib:KondKopLek}, Л.~М.~Кожевниковой \cite{KozhKhad:bib:Kozhevnikova},

Ф.~Х.~Мукминова, В.~Ф.~Гилимшиной \cite{KozhKhad:bib:GilMuk} и других авторов.

В работе \cite{KozhKhad:bib:KozhKarimov} Л.~М.~Кожевниковой, Р.~Х.~Каримовым для решений   квазилинейных эллиптических уравнений второго порядка установлены оценки  сверху и доказана их точность.

Л.~М.~Кожевниковой, А.~А.~Хаджи \cite{KozhKhad:bib:KozhKhad1} этот результат обобщён на некоторый класс анизотропных эллиптических уравнений второго порядка.

Для уравнений с нестепенными нелинейностями исследование скорости убывания решений на бесконечности в неограниченных областях ранее не проводилось.



\smallskip



\Section{ Пространства Соболева"--~Орлича}

Приведём необходимые сведения из теории $N$-функций и пространств Соболева"--~Орлича \cite{KozhKhad:bib:Rut}.

Неотрицательная непрерывная выпуклая вниз функция $M(z)$, $z\in\mathbb{R}_{1}$ называется $N$-функцией, если она

чётна и

$$\lim\limits_{z \to  0}\frac{M(z)}{z}=0, \quad

\lim\limits_{z \to \infty}\frac{M(z)}{z}=\infty.$$

Отметим, что $M(\varepsilon u)\leq \varepsilon M(u)$ при $0<\varepsilon\leq1$.

Для $N$-функции $M(z)$ имеет место интегральное представление

$$M(z)=\int_{0}^{|z|}m(t)dt,$$ где $m(t)$ "--- положительная при $t>0$, не убывающая  и непрерывная справа функция при $t\geq0$ такая, что $$m(0)=0,\quad \lim\limits_{t \to  \infty}m(t)=\infty.$$



$N$-функция

$$

\overline{M}(z)=\sup_{y\geq 0}(y|z|-M(y))

$$

называется дополнительной к $N$-функции $M(z)$. Известно следующее неравенство Юнга:

\begin{equation}\label{eq:Kozh:ung}

|zy|\leq M(z)+\overline{M}(y),\quad z,y\in \mathbb{R}.

\end{equation}

Кроме того, имеет место равенство

\begin{equation}\label{eq:Kozh:ung-r}

|z|m(|z|)=\overline{M}(m(|z|))+M(z),\quad  z\geq z_0 \quad  z\in \mathbb{R}.

\end{equation}



Для $N$-функций $P(z)$, $M(z)$ записывают $P(z)\prec M(z)$, если существуют числа $l>0$, $z_0\geq 0$ такие, что

$$

P(z)\leq M(lz).

$$

$N$-функции $P(z)$, $M(z)$ называются сравнимыми, если имеет место одно из соотношений: $P(z)\prec M(z)$ или $M(z)\prec P(z)$.

$N$-функции $P(z)$ и $M(z)$ называются эквивалентными, если $P(z)\prec M(z)$ и $M(z)\prec P(z)$.



$N$-функция $M(z)$ удовлетворяет $\Delta_2$-условию  при больших значениях $z$, если существуют

такие числа $c>0$, $z_0\geq0$, что $M(2z)\leq cM(z)$ для любых $z\geq z_0$.

$\Delta_2$-условие  эквивалентно выполнению  неравенства

\begin{equation}\label{eq:Kozh:delta2}

M(lz)\leq c(l)M(z),\quad z\geq z_0,

\end{equation}

где $l$ "--- любое число, большее единицы, $c(l)>0$.



$N$-функция $M(z)$ удовлетворяет $\Delta_2$-условию

тогда и только тогда, когда существуют положительные числа $c>1$, $z_0\geq 0$ такие, что при $z\geq z_0$

справедливо неравенство

\begin{equation}\label{eq:Kozh:usl}

  {zm(z)}<c{M(z)}.

\end{equation}

В каждом классе эквивалентных $N$-функций, подчиняющихся $\Delta_2$-условию, имеются $N$-функции, удовлетворяющие неравенству \eqref{eq:Kozh:delta2} при всех $z\geq 0$.

В дальнейшем в работе предполагается, что $\Delta_2$-условие для рассматриваемых $N$-функций выполняется  при всех значениях $z\geq0$  (т.~е. $z_0=0$).



Для $N$-функции $M(z)$ ввиду выпуклости и неравенства \eqref{eq:Kozh:delta2} существует $c>0$ такое, что справедливо неравенство

\begin{equation}\label{eq:Kozh:sum}

M(y+z)\leq c M(z)+c M(y),\quad z,y\geq 0.

\end{equation}



Пусть $Q\subset\mathbb{R}_n$. Классом Орлича $K_M(Q)$, соответствующем $N$-функции $M(z)$, называется множество измеримых в $Q$

функций $v$ таких, что

$$

\int_Q M(v({\bf x}))d{\bf x}<\infty.

 $$

Пространством Орлича $L_M(Q)$ называется линейная оболочка $K_M(Q)$  с нормой Люксембурга

$$

 \|v\|_{L_{M}(Q)}=\|v\|_{M,Q}=\inf\left\{k\geq 0 \;\Big | \; \int_{Q}M\Bigl(\frac{v(\mathbf{x})}{k}\Bigr) d{\mathbf{ x}}\leq1\right\}.

$$

При $Q=\Omega$ будем использовать обозначение $\|\cdot\|_{M,\Omega}=\|\cdot\|_M$.  Класс Орлича $K_M(Q)$ совпадает с

пространством Орлича $L_M(Q)$ тогда и только тогда, когда $M(z)$ удовлетворяет $\Delta_2$-условию.



Для функций $v\in L_{M}(Q)$, $u\in {L}_{\overline {M}}(Q)$ имеют место неравенство Гёльдера

 \begin{equation}\label{eq:Kozh:gel}

\left | \int_Q u({\bf x})v({\bf x})d{\bf x}\right|\leq 2\|v\|_{M,Q}\|u\|_{\overline {M},Q}

\end{equation}

и  неравенство

\begin{equation}\label{eq:Kozh:ner_1}

\|v\|_{M,Q}\leq\int_{Q}M(v)d{\bf x}+1.

\end{equation}

Кроме того, если $N$-функция $M(z)$ удовлетворяет $\Delta_2$-условию, то для $v\in L_M(Q)$

справедливы неравенства

\begin{multline}\label{eq:Kozh:ner_2}

 \int_Q M(v)d{\bf x}=\int_Q M\left(\|v\|_{M,Q}\frac{v}{\|v\|_{M,Q}}\right)d{\bf

 x}\leq\int_Q M\left(l\frac{v}{\|v\|_{M,Q}}\right)d{\bf x}\leq

\\

\leq

 c(l)\int_Q M\left(\frac{v}{\|v\|_{M,Q}}\right)d{\bf x}= c(l),\quad

 l=\max\left(\|v\|_{M,Q}, 1\right).

\end{multline}



\smallskip

\hypertarget{kozh:lemma1}{}

\begin{lemma}[1]Если $N$-функция $M(z)$ удовлетворяет $\Delta_2$-условию,  $v({\bf x}),$ $v^i({\bf x})\in L_M(Q),$ $i=1, 2,\ldots,$ $v^i({\bf x}) \to  v({\bf x})$ в $L_M(Q),$ то

\begin{equation}\label{eq:Kozh:lem_1}

\int_Q M(v^i)d{\bf x} \to \int_Q M(v)d{\bf x},\quad i \to  \infty.

\end{equation}

\end{lemma}







\begin{proof} Ввиду справедливости $\Delta_2$-условия сходимость по норме равносильна сходимости в среднем \cite[гл. II, \S~9, теорема 9.4]{KozhKhad:bib:Rut}, следовательно, для любого $\varepsilon>0$ существует номер $i_0$ такой, что

\begin{equation}\label{eq:Kozh:lem_1(1)}

\int_Q M(v^i-v)d{\bf x}< \varepsilon,\quad i\geq i_0.

\end{equation}

 Сходимость $v^i({\bf x}) \to  v({\bf x})$ в $L_M(Q)$  влечёт ограниченность множества $\|v^i\|_{M,Q}$, $i=1, 2,\ldots$. В свою очередь, из

неравенства  \eqref{eq:Kozh:ner_2} следует существование числа $C_1>0$ такого, что

\begin{equation}\label{eq:Kozh:lem_1(2)}

M(v^i)\leq C_1,\quad i=1, 2,\ldots.

\end{equation}







Далее, пользуясь формулой Лагранжа, неравенствами \eqref{eq:Kozh:ung}, \eqref{eq:Kozh:delta2}, для любого $\varepsilon\in (0, 1]$ и  фиксированного $\theta\in [0, 1]$ выводим

\begin{multline}\label{eq:Kozh:lem_1(3)}

\left|\int_Q

(M(v^i)-M(v)) d{\bf x}\right|\le\int_Q |v^i-v|m(|\theta v^i+(1-\theta)v|)

d{\bf x}\leq

\\

\le\int_Q \left\{M\Bigl(\frac{1}{\varepsilon}(v^i-v)\Bigr)+\overline{M}\left(\varepsilon

m(|\theta v^i+(1-\theta)v|)\right)\right\} d{\bf x}\leq

\\

\leq \int_Q \left\{

C_2M(v^i-v)+\varepsilon\overline{M}\left(

m(|\theta v^i+(1-\theta)v|)\right)\right\} d{\bf x}.

\end{multline}

Применяя \eqref{eq:Kozh:ung-r}, \eqref{eq:Kozh:usl}, \eqref{eq:Kozh:sum}, устанавливаем цепочку неравенств

\begin{multline}\label{eq:Kozh:lem_1(4)}

\overline{M}\left(

m(|\theta v^i+(1-\theta)v|)\right)\leq

|\theta v^i+(1-\theta)v|m(|\theta v^i+(1-\theta)v|)\le

\\

\leq c M(\theta v^i+(1-\theta)v)\leq C_3(M(v^i)+M(v)).

\end{multline}



Соединяя \eqref{eq:Kozh:lem_1(3)}, \eqref{eq:Kozh:lem_1(4)}, \eqref{eq:Kozh:lem_1(2)}, \eqref{eq:Kozh:lem_1(1)}, при $i\geq i_0$ получаем неравенства

$$\left|\int_Q

\left(M(v^i)-M(v)\right) d{\bf x}\right|\leq C_2 \int_Q

M(v^i-v)d{\bf x}+\varepsilon C_3\int_Q(M(v^i)+M(v))d{\bf x}\leq C_4 \varepsilon.

$$

Таким образом, сходимость \eqref{eq:Kozh:lem_1} доказана.

\end{proof}



\smallskip





Пусть $B_1(z),\dots,B_n(z)$ "--- $N$-функции, определим пространство Соболева"--~Орлича $\stackrel

{\circ}

{H}\!_{B}^1(Q)$ как пополнение $C_0^{\infty}(Q)$ по

норме

$$

\|v\|_{\stackrel {\circ}

{H}{}_{B}^1(Q)}=\sum_{\alpha=1}^n\|v_{x_\alpha}\|_{B_\alpha,Q}.$$

Положим

$$

h(t)=t^{-\frac{1}{n}}\biggl(\prod_{\alpha=1}^n B_{\alpha}^{-1}(t)\biggr)^{{1}/{n}}

$$

и будем предполагать, что

$$\int_0^1{t}^{-1}{h(t)}dt<\infty.$$

Тогда можно определить $N$-функцию $B^{*}(z)$ по формуле

$$

(B^{*})^{-1}(z)=\int_0^{|z|} {t}^{-1}{h(t)}dt.

$$

Теоремы вложения для пространства $\stackrel {\circ}

{H}\!_{B}^1(Q)$ установлены в \cite{KozhKhad:bib:Klimov3, KozhKhad:bib:Korolev1}. Приведем теорему вложения  А.~Г.~Королёва~\cite{KozhKhad:bib:Korolev1}, доказанную для ограниченных областей $Q$.



\smallskip



\hypertarget{kozh:lemma2}{}



\begin{lemma}[2]Пусть $v\in \stackrel {\circ} {H}\!_{B}^1(Q).$

\begin{itemize}



\item[\rm 1)] Если

\begin{equation}\label{eq:Kozh:1}

\int_1^{\infty}t^{-1}{h(t)}dt=\infty,

\end{equation}

то

\begin{equation}\label{eq:Kozh:K1}

\|v\|_{B^*,Q}\leq A_1\sum_{\alpha=1}^n\|v_{x_{\alpha}}\|_{B_{\alpha},Q};

\end{equation}

\item[\rm 2)] если

\begin{equation}\label{eq:Kozh:2}

\int_1^{\infty}t^{-1}{h(t)}dt<\infty,

\end{equation}

то

\begin{equation}\label{eq:Kozh:K2}

\sup_{Q}|v|\leq A_2\sum_{\alpha=1}^n\|v_{x_{\alpha}}\|_{B_{\alpha},Q}.

\end{equation}



\end{itemize}

Здесь $$A_1=\frac{n-1}{n},\quad A_2=\int_0^\infty\frac{h(t)}{t}dt.$$

\end{lemma}





\smallskip



\begin{remark}[1]

 Лемма \hyperlink{kozh:lemma2}{2} справедлива также и для произвольных не\-ог\-раниченных областей $\Omega\subset\mathbb{R}_n$. Действительно, пусть $v({\bf x})\in\stackrel {\circ}{H}\!_{B}^1(\Omega)$, тогда существует последовательность $v^i({\bf x})\in C_0^\infty(\Omega)$ такая, что $v^i({\bf x}) \to  v({\bf x})$ в $\stackrel {\circ}{H}\!_{B}^1(\Omega)$. Записывая неравенства \eqref{eq:Kozh:K1}, \eqref{eq:Kozh:K2} для функций $v^i$ и выполняя предельный переход при $i \to \infty$, установим их и для функции $v\in\stackrel {\circ}{H}\!_{B}^1(\Omega)$.

\end{remark}



\smallskip



\begin{example}[1]Возьмём $n=2$, $B_1(z)=z^{[{5}/{4}, {3}/{2}]}$, $B_2(z)=z^{[{5}/{3}, 2]}$.

Здесь использовано обозначение

$$

z^{[a,b]}=

\left \{

\begin{array}{ll}

z^a, \  \mbox{  } \ z\leq 1,\\

z^b, \  \mbox{  } \ z> 1.\\

\end{array}\right.

$$

Поскольку

$$h(t)=t^{[{1}/{5},{1}/{12}]},\quad \int_0^1t^{-1}{h(t)}dt<\infty,\quad \int_1^{\infty}t^{-1}{h(t)}dt=\infty,$$

 можно определить функции

$$(B^{*})^{-1}(z)=\int_0^{|z|}\frac{t^{[{1}/{5}, {1}/{12}]}}{t} dt=

\left \{

\begin{array}{ll}

5|z|^{{1}/{5}}, & z\leq1\\

12|z|^{{1}/{12}}-7, & z>1,\\

\end{array}\right.$$

$$

B^{*}(z)=

\left\{

\begin{array}{ll}

(z/5)^5, & z\leq 5,\\

((z+7)/12)^{12}, & z>5.\\

\end{array}\right.

$$

\end{example}

При этом справедливо неравенство \eqref{eq:Kozh:K1}.



\smallskip



\begin{example}[2]Положим $n=3$, $B_1(z)=z^{[{3}/{2}, {7}/{2}]}$, $B_2(z)=z^{[3, 7]}$, $B_3(z)=z^{[3, {7}/{3}]}$.

Поскольку

$$h(t)=t^{[{1}/{9},- {1}/{21}]},\quad \int_0^\infty t^{-1}{h(t)}dt<\infty,$$ справедливо неравенство \eqref{eq:Kozh:K2}.

\end{example}



\smallskip



Если $N$-функция $B^*(z)$ удовлетворяет $\Delta_2$-условию, то из \eqref{eq:Kozh:delta2} следует, что определена и конечна функция

$$

\Lambda(z)=\sup\limits_{t>0}\frac{B^*(zt)}{B^*(t)}

$$

и при всех значениях $t,z\in \mathbb{R}$ справедливы  неравенства

$$B^*(tz)\leq B^*(t)\Lambda(z),\quad \Lambda(zt)\leq \Lambda(z)\Lambda(t).$$

Приведём еще одну теорему вложения

 А.~Г.~Королёва~\cite{KozhKhad:bib:Korolev2}, доказанную также для ограниченных областей $Q$.



\smallskip

\hypertarget{kozh:lemma3}{}

\begin{lemma}[3]Пусть $v\in \stackrel {\circ}

{H}{}_{B}^1(Q),$ выполнено условие \eqref{eq:Kozh:1} и функции $B_{\alpha}(z),$ $\alpha=1, 2,\ldots,n,$ $B^*(z)$ удовлетворяют $\Delta_2$-условию,  тогда справедливо неравенство

\begin{equation}\label{eq:Kozh:K3}

\int_{Q}B^*(v({\bf x}))d{\bf x}\leq \chi\left(A_3\sum_{\alpha=1}^n \int_{Q}B_{\alpha}(v_{x_{\alpha}})d{\bf x}\right).

\end{equation}

Здесь $A_3$ зависит от $n,$ $\chi(z)=z\Lambda(z^{1/n}),$ $z\geq 0.$

\end{lemma}



\smallskip





\begin{remark}[2]

Лемма \hyperlink{kozh:lemma3}{3} справедлива также и для произвольных не\-ог\-ра\-ниченных областей $\Omega\subset\mathbb{R}_n$.

Действительно, пусть $v({\bf x})\in\stackrel {\circ}{H}\!_{B}^1(\Omega)$, тогда существует последовательность $v^i({\bf x})\in C_0^\infty(\Omega)$ такая, что $v^i({\bf x}) \to  v({\bf x})$ в $\stackrel {\circ}{H}\!_{B}^1(\Omega)$. Применяя лемму \hyperlink{kozh:lemma1}{1}, имеем

$$\int_{\Omega}B_{\alpha}(u^i_{x_{\alpha}})d{\bf x} \to \int_{\Omega}B_{\alpha}(u_{x_{\alpha}})d{\bf x},\quad \alpha=1,\ldots,n,$$

$$ \int_{\Omega}B^*(u^i)d{\bf x} \to  \int_{\Omega}B^*(u)d{\bf x},\quad i \to \infty.$$

  Записывая неравенство \eqref{eq:Kozh:K3} для функций $u^i$, пользуясь непрерывностью    функции $\chi(z)$,  выполняя предельный переход при $i \to \infty$, устанавливаем его и для функции $u\in\stackrel {\circ}{H}\!_{B}^1(\Omega)$.

\end{remark}



\smallskip



\Section{Формулировка основных результатов}

Пусть  существуют неотрицательные функции $\psi_1({\bf x}),\psi({\bf x})\in L_1(\Omega)$, $\varphi^{-1}({\bf x})$, $\varphi({\bf x})$, $\varphi_1({\bf x})\in L_{\infty}(\Omega)$ такие, что  для п.~в.  ${\bf x}\in\Omega$, ${\bf p}=(p_0,p)$, ${\bf q}=(q_0,q)\in\mathbb{R}_{n+1}$, ${\bf p}\neq {\bf q}$, справедливы следующие неравенства:

\begin{gather}

\label{eq:Kozh:us1}

\sum_{\alpha=0}^na_{\alpha}({\bf x},p_0,p)p_\alpha\geq \varphi({\bf x})\sum_{\alpha=0}^nB_\alpha(p_\alpha)-\psi({\bf x});

\\

\label{eq:Kozh:us2}

 \sum_{\alpha=0}^n\overline{B}_\alpha(a_{\alpha}({\bf x},p_0,p))\leq \varphi_1({\bf x}) \sum_{\alpha=0}^n

B_\alpha(p_\alpha)+\psi_1({\bf x});

\\

\label{eq:Kozh:us4}

 \sum_{\alpha=0}^n (a_\alpha({\bf x},p_0,p)-a_\alpha({\bf x},q_0,q))(p_{\alpha}-q_{\alpha})> 0.

\end{gather}

Здесь $B_0(z),B_1(z),\dots,B_n(z)$ "--- $N$--функции, удовлетворяющие $\Delta_2$-условию. В случае выполнения условия \eqref{eq:Kozh:1} будем считать

$$

B_0(z)\prec B^*(z),

$$

а при выполнении \eqref{eq:Kozh:2} $B_0(z)$ "--- произвольная $N$-функция.



  Определим пространство Соболева"--~Орлича $\stackrel {\circ}

{W}\!_{\bf B}^{1}(\Omega)$ как пополнение $C_0^{\infty}(\Omega)$ по

норме

$$

\|u\|_{\stackrel {\circ}

{W}{}_{\bf B}^{1}(\Omega)}=\|u\|_{B_0}+\|u\|_{\stackrel {\circ}

{H}{}_{B}^1(\Omega)}.

$$

Определим оператор  ${\bf B}$ : $\stackrel {\circ}

{W}{}_{\bf B}^{1}(\Omega) \to  L_1(\Omega)$  формулой

$$

{\bf B}(v)=B_0(v)+\sum_{\alpha=1}^n B_\alpha(v_{x_{\alpha}}),\quad v\in\stackrel {\circ}

{W}{}_{\bf B}^{1}(\Omega).

$$

Пусть $Q\subseteq \Omega$ ($Q$ может совпадать c $\Omega$), для $v\in \stackrel {\circ}{W}\!\!_{{\bf B}}^{1}(\Omega)$ определим функционал ${\bf A}(u)$ $\bigl(u\in \stackrel {\circ}{W}\!\!_{{\bf B}}^{1}(\Omega)\bigr)$

  равенством

\begin{equation}

\label{eq:Kozh:A}

({\bf A}(u),v)_Q=\int_{Q}\biggl(\sum\limits_{\alpha=1}^{n}a_{\alpha}({\bf x},u,\nabla u)v_{x_{\alpha}}+a_0({\bf x},u,\nabla u)v\biggr)d{\bf x}.

\end{equation}

Если $Q=\Omega$, то будем писать $({\bf A}(u),v)_{\Omega}=({\bf A}(u),v)$.

Через $\|\cdot\|_{m,Q}$ будем обозначать норму в пространстве $L_m(Q),\;$ $1\leq m\leq \infty$, при  $Q=\Omega$ индекс $Q$ будет опускаться.



\smallskip

\begin{definition}[1]\it Обобщённым решением задачи \rm \eqref{eq:Kozh:ur}, \eqref{eq:Kozh:gu}  назовём функцию $u({\bf x}) \in\stackrel {\circ} {W}\!{}^{1}_{\mathbf{B}} ({\Omega})$, удовлетворяющую

интегральному тождеству

\begin{equation}\label{eq:Kozh:intt}

({\bf A}(u),v)=0

\end{equation}

для любой функции $v({\bf x}) \in\stackrel {\circ} {W}\!{}^{1}_{\bf B} (\Omega)$.

\end{definition}





\smallskip





Пусть $N$-функции

$$B_{\alpha}(z)=\int_0^{|z|}b_{\alpha}(t)dt$$ кроме $\Delta_2$-условия  удовлетворяют требованию

\begin{equation}\label{eq:Kozh:dop_usl}

\lim_{\lambda\to\infty}\inf\limits_{z> 0}\frac{b_\alpha(\lambda z)}{b_{\alpha}(z)}=\infty, \quad \alpha=0, 1,\ldots,n.

\end{equation}



\smallskip

\begin{example}[3]Легко проверить, что функция

$$B(u)=|u|^{a}(|\ln|u||+1),\quad a>1,$$

удовлетворяет $\Delta_2$-условию  и \eqref{eq:Kozh:dop_usl}.

\end{example}



\smallskip



\hypertarget{kozh:th1}{}

\begin{theorem}[1] Пусть выполнены условия \eqref{eq:Kozh:us1}--\eqref{eq:Kozh:us4}{\rm ,} \eqref{eq:Kozh:dop_usl}{\rm ,} тогда существует единственное решение $u({\bf x})$ задачи \eqref{eq:Kozh:ur}{\rm ,}

\eqref{eq:Kozh:gu}   и $\mathcal{M}_1>0$ такое{\rm ,} что  справедлива оценка

\begin{equation}\label{eq:Kozh:the_3}

\|\mathbf{B}(u)\|_1\leq

\mathcal{M}_1\|\psi\|_1.

\end{equation}

\end{theorem}



\smallskip



При дополнительном требовании, согласно которому

\begin{equation}\label{eq:Kozh:us5}

\mbox {функция}\;  \sum \limits_{\alpha=1}^n a_{\alpha}({\bf x},p_0,p)p_{\alpha}

\left[

\begin{array}{ll}

\mbox {монотонно не убывает}&  \mbox {при} \quad p_0\geq 0,\\

\mbox {монотонно не возрастает}&  \mbox  {при}\quad  p_0<0

\end{array}

\right],

\end{equation}

доказана ограниченность решения задачи \eqref{eq:Kozh:ur}, \eqref{eq:Kozh:gu}.



\smallskip



\hypertarget{kozh:th2}{}



\begin{theorem}[2]Пусть выполнены условия  \eqref{eq:Kozh:us1}--\eqref{eq:Kozh:us4}{\rm ,} \eqref{eq:Kozh:us5}{\rm ,} \eqref{eq:Kozh:dop_usl}{\rm ,} \eqref{eq:Kozh:1}{\rm ,} $N$"=функция $B^*(z)$ удовлетворяет $\Delta_2$-условию   и

\begin{equation}\label{eq:Kozh:the_2_1}

\psi\in L_{l/(l-1)}(\Omega),\quad l\geq 1,

\end{equation}

где $l$ такое{\rm ,} что $t^{1/l-1}\Lambda(t^{1/(ln)}) \to  0$ при $t \to  0.$ Тогда для обобщённого решения задачи \eqref{eq:Kozh:ur}{\rm ,} \eqref{eq:Kozh:gu} $u({\bf x})$ существует $\mathcal{M}_2>0$ такое{\rm ,} что справедлива оценка

\begin{equation}\label{eq:Kozh:the_2_2}

\sup\limits_{\Omega}|u|\leq \mathcal{M}_2.

\end{equation}

\end{theorem}



\smallskip



\begin{remark}[3]\label{mark_4}

В случае выполнения условия \eqref{eq:Kozh:2} ограниченность решения задачи \eqref{eq:Kozh:ur}, \eqref{eq:Kozh:gu} следует из неравенств \eqref{eq:Kozh:K2}, \eqref{eq:Kozh:the_3}.

\end{remark}



\smallskip



Далее приведём результат для  областей, расположенных вдоль выделенной оси

$Ox_s,$ $s\in\{1, 2,\ldots,n\}$ (область $\Omega$ лежит в

полупространстве $x_s>0$, сечение $\gamma_r=\{\textbf{x}\in\Omega

\;|\; x_s=r\}$ не пусто, связно и ограничено при любом $r>0$). Ниже будет

использовано следующее обозначение: $$\Omega_{c}^{d}[s]=\{{\bf

x}\in\Omega\;|\;c<x_s<d\},$$ при этом значения $c=0,\;d=\infty$

опускаются.



C целью изучения поведения решения задачи \eqref{eq:Kozh:ur}, \eqref{eq:Kozh:gu} при

$x_s \to  \infty$ определим функции

\begin{equation}\label{eq:Kozh:nu}

\nu_i (r)=\inf_{v\in C_0^{\infty}(\Omega)}\sup\left\{z\;\Big|\; \int_{\gamma_r}B_i(zv)d{\bf x}'_s\leq\int_{\gamma_r}B_i(v_{x_i})d{\bf x}'_s\right\},\; \end{equation}

$$\overline{\nu}_s(x)=\min\limits_{i=\overline{1,n},i\neq s}\nu_i(x),$$

где ${\bf x}'_s=\{x_1,x_2,\dots,x_{s-1},x_{s+1},\dots,x_n\}$.



 %будем считать, что область $\Omega$

%удовлетворяет условию

 %\begin{eqnarray}\label{infty}

%\int\limits_1^{\infty}\overline{\nu}_s(r)dr=\infty.

%\end{eqnarray}



Пусть существуют числа $\beta_i\geq 0$, $i=1,\ldots,n$, $i\neq s$, $0<z^*\leq \infty$ такие, что

\begin{equation}\label{eq:Kozh:BB}

B_s(z)\leq\sum\limits_{i=1,i\neq s}^n \beta_i B_i(z),\quad 0\leq z\leq z^*.

\end{equation}

Предположим, что

\begin{equation}\label{eq:Kozh:supp}

 {\rm supp}\;  \psi_1,\psi\subset\{{\bf x}\in \mathbb{R}_n \;|\;x_s<R_0\},\quad R_0>0.

\end{equation}



В следующей теореме установлена оценка, характеризующая скорость убывания решения при $x_s \to  \infty$.



\smallskip



\hypertarget{kozh:th3}{}

\begin{theorem}[3]Пусть область $\Omega$

расположена вдоль оси $Ox_s,$ $s\in\{1, 2,\ldots,n\}$ и выполнены условия  \eqref{eq:Kozh:us1}--\eqref{eq:Kozh:us4}{\rm ,} \eqref{eq:Kozh:dop_usl}{\rm ,} \eqref{eq:Kozh:BB}{\rm ,} \eqref{eq:Kozh:supp}{\rm .} Если $z^*\neq \infty,$ то решение задачи \eqref{eq:Kozh:ur}--\eqref{eq:Kozh:gu} предполагается ограниченным{\rm .}

Тогда существуют положительные числа $\kappa,$

$\mathcal{M}$ такие{\rm ,} что

 при всех  $r\geq 2R_0$

справедлива  оценка

\begin{equation}\label{eq:Kozh:the_4}

\|\mathbf{B}(u)\|_{1,\Omega_r}\leq \mathcal{M}

\exp\biggl(-\kappa\int_{2R_0}^r \overline{\nu}_s^*(\rho)d\rho\biggr),

\end{equation}

где $\overline{\nu}_s^*(\rho)=\min\{\overline\nu_s(\rho),z^*\}.$

\end{theorem}



\smallskip



\Section{Существование и единственность обобщённого решения}

Обозначим \linebreak $\|\varphi\|_{\infty}=\widehat{a}$, $\|1/\varphi\|_{\infty}=1/\overline{a}$, $\|\varphi_1\|_{\infty}=\widehat{a}_1$.

Покажем, что определение обобщённого решения задачи \eqref{eq:Kozh:ur}, \eqref{eq:Kozh:gu} корректно.

 Из условия  \eqref{eq:Kozh:us2} и неравенства \eqref{eq:Kozh:ner_1} следует оценка

\begin{multline}\label{eq:Kozh:oc_a_1}

\sum_{\alpha=0}^n\|a_\alpha({\bf x},u,\nabla u)\|_{\overline{B}_{\alpha}}\leq

\\

\leq\sum_{\alpha=0}^n\int_{\Omega}\overline{B}_{\alpha}(a_\alpha({\bf x},u,\nabla u))d{\bf x}+n+1\leq \widehat{a}_1\|\mathbf{B}(u)\|_{1}+\|\psi_1\|_{1}+n+1.

\end{multline}



Используя неравенство Гёльдера \eqref{eq:Kozh:gel} и оценку \eqref{eq:Kozh:oc_a_1}, для функций $u(\mathbf{x})$, $v(\mathbf{x})\in$ $

\stackrel {\circ} {W}\!{}^{1}_{\mathbf{B}} (\Omega)$

выводим

\begin{multline}\label{eq:Kozh:opr1_1}

|(\mathbf{A}(u),v)|\leq\int_{\Omega}

\biggl(\sum_{\alpha=1}^n|a_\alpha({\bf x},u,\nabla

u)||v_{x_\alpha}|+|a_0({\bf x},u,\nabla u)||v|\biggr)d{\bf x}\leq

\\

\leq2\sum_{\alpha=1}^n\|a_\alpha({\bf x},u,\nabla u)\|_{\overline{B}_{\alpha}}\|v_{x_\alpha}\|_{{B}_{\alpha}}+2\|a_0({\bf x},u,\nabla u)\|_{\overline{B}_0}\|v\|_{{B}_0}\leq \\ \leq C_1\|v\|_{\stackrel {\circ}

{W}{}_{\bf B}^{1}(\Omega)}<\infty.

\end{multline}

Таким образом, интегралы, входящие в \eqref{eq:Kozh:intt},  конечны.



\smallskip



\begin{remark}[4]

Из условий \eqref{eq:Kozh:dop_usl} следует существование $\lambda_0>0$ такого, что для любых $\lambda\geq\lambda_0$ и любых $z\geq 0$ справедливы неравенства

$$b_{\alpha}(\lambda z)\geq 2 b_{\alpha}(z),\quad \alpha=0, 1,\ldots,n.$$

Последние неравенства обеспечивают выполнение для любых $z\in\mathbb{R}$ условий

$$B_{\alpha}(z)\leq \frac{1}{2\lambda}B_{\alpha}(\lambda z),\quad \lambda\geq\lambda_0,\quad \alpha=0, 1,\ldots,n.

$$

 Тогда, согласно  \cite[I, \S~4, теорема 4.2]{KozhKhad:bib:Rut}, $N$-функции $\overline{B}_{\alpha}(z)$, $\alpha=0, 1,\ldots,n$, удовлетворяют $\Delta_2$-условию.

\end{remark}



\smallskip



{\it Д\,о\,к\,а\,з\,а\,т\,е\,л\,ь\,с\,т\,в\,о \ т\,е\,о\,р\,е\,м\,ы \rm \hyperlink{kozh:th1}{1}.}

Единственность следует из условия строгой монотонности \eqref{eq:Kozh:us4}.

Полагая  в  равенстве \eqref{eq:Kozh:intt} $v=u$,  применяя \eqref{eq:Kozh:us1}, получаем неравенство \eqref{eq:Kozh:the_3}.

 Точно так же выводится неравенство

\begin{equation}\label{eq:Kozh:the_3(4)}

\overline{a}\|\mathbf{B}(u)\|_{1}\leq \|\psi\|_{1}+(\mathbf{A}(u),u).

\end{equation}





По элементу $u\in\stackrel {\circ}

{W}\!_{\bf B}^{1}(\Omega)$ определим элемент $\mathbf{A}(u)\in\bigl(\stackrel {\circ}

{W}\!_{\bf B}^{1}(\Omega)\bigr)'$ равенством \eqref{eq:Kozh:A}.

Оператор $\mathbf{A}$ определён (см. \eqref{eq:Kozh:opr1_1}), проверим условия теоремы Ж.~Л.~Лионса \cite[гл. II, \S~2, теорема 2.1]{KozhKhad:bib:Lions}.



\smallskip



1) Из неравенств \eqref{eq:Kozh:opr1_1}, \eqref{eq:Kozh:oc_a_1}, согласно \eqref{eq:Kozh:ner_1}, следует слабая ограниченность множества $\{\mathbf{A}(u), u\in \Theta\}$ на ограниченном множестве $\Theta\subset\stackrel {\circ}

{W}\!_{\bf B}^{1}(\Omega)$. Тогда $\{\mathbf{A}(u), u\in \Theta\}$ ограничено в $\left(\stackrel {\circ}

{W}\!_{\bf B}^{1}(\Omega)\right)'$, т.~е. оператор $\mathbf{A}$ ограничен.



\smallskip



2) Монотонность оператора $\mathbf{A}$ обеспечивается условием \eqref{eq:Kozh:us4}.



\smallskip



3) Докажем коэрцитивность оператора $\mathbf{A}$. Пользуясь \eqref{eq:Kozh:the_3(4)}, выводим

\begin{multline}\label{eq:Kozh:the_3(5)}

\frac{(\mathbf{A}(u),u)}{\|u\|_{

\stackrel {\circ}

{W}\!_{\bf B}^{1}(\Omega)}}\geq \frac{1}{\|u\|_{

\stackrel {\circ}

{W}\!_{\bf B}^{1}(\Omega)}}\left(\overline{a}\|\mathbf{B}(u)\|_{1}-\|\psi\|_{1}\right)=

\\

=\frac{1}{\|u\|_{\stackrel {\circ}

{H}{}_{\bf B}^1(\Omega)}+\|u\|_{B_0}}\left({\overline{a}}\int_{\Omega}\biggl(\sum_{\alpha=1}^nB_\alpha(u_{x_\alpha})+B_0(u)\biggr)d{\bf x}-C_2\right).

\end{multline}

Далее из условий \eqref{eq:Kozh:dop_usl} следует, что для любого $R>0$ найдётся $\lambda_0>0$ такое, что при

достаточно больших $\|u\|_{B_0}>\lambda_0$, $\|u_{x_{\alpha}}\|_{B_{\alpha}}>\lambda_0$, $\alpha={1,\ldots,n}$,

 справедливы неравенства

\begin{equation}\label{eq:Kozh:the_3(6)}

b_0(|u|)>Rb_0\Bigl(\frac{|u|}{\|u\|_{B_0}}\Bigr),\quad

b_{\alpha}(|u_{x_{\alpha}}|)>Rb_{\alpha}\Bigl(\frac{|u_{x_{\alpha}}|}{\|u_{x_{\alpha}}\|_{B_{\alpha}}}\Bigr),\;\alpha={1,\ldots,n}.

\end{equation}



Пусть $\|u^i\|_{

\stackrel {\circ}

{W}\!_{\bf B}^{1}(\Omega)} \to \infty$ при $i \to \infty$, можно

считать, что

$$\|u^{i}_{x_1}\|_{B_1}+\ldots+\|u^{i}_{x_n}\|_{B_n}+\|u^{i}\|_{B_0}>(n+1)\lambda_0,\quad

i\geq i_0.$$

При каждом $i\geq i_0$ найдётся хотя бы одно слагаемое больше $\lambda_0$. Пусть для определённости при фиксированном  $i\geq i_0$ наибольшим является

первое слагаемое: $$\|u^i_{x_1}\|_{B_1}>\lambda_0.$$





Соединяя \eqref{eq:Kozh:the_3(5)},  \eqref{eq:Kozh:usl}, получаем

$$

\frac{(\mathbf{A}(u^{i}),u^{i})}{\|u^{i}\|_{

\stackrel {\circ}

{W}\!_{\bf B}^{1}(\Omega)}}\geq\frac{\overline{a}}{c}\frac{1}{(n+1)\|u^{i}_{x_1}\|_{B_1}}

\int_{\Omega}|u^{i}_{x_1}|b_1(|u^{i}_{x_1}|)d{\bf

x}-\frac{C_2}{(n+1)\lambda_0}.

$$

Далее, применяя \eqref{eq:Kozh:the_3(6)}, выводим

\begin{multline*}

\frac{(\mathbf{A}(u^{i}),u^{i})}{\|u^{k}\|_{

\stackrel {\circ}

{W}\!_{\bf B}^{1}(\Omega)}}\geq R\frac{C_3}{\|u^{i}_{x_1}\|_{B_1}}

\int_{\Omega}|u^{i}_{x_1}|b_1\Bigl(\frac{|u^{i}_{x_1}|}{\|u^{i}_{x_1}\|_{B_1}}\Bigr)d{\bf

x}-{C_4}\geq

\\

\geq R{C_3}

\int_Q

B_1\Bigl(\frac{u^{i}_{x_1}}{\|u^{i}_{x_1}\|_{B_1}}\Bigr)d{\bf

x}-C_4=

RC_3-C_4.

\end{multline*}





Таким образом, ввиду произвольности $R$ и $i$, $i\geq i_0$, установлена коэрцетивность оператора $\mathbf{A}$.



\smallskip



4) Семинепрерывность оператора $\mathbf{A}$ вытекает из непрерывности функций $a_{\alpha}({\bf x},p_0,p)$, $\alpha=0, 1,\ldots,n$, по $(p_0,p)\in \mathbb{R}_n$ и теоремы Лебега.



 Согласно  теореме Ж.~Л.~Лионса \cite[гл. II, \S~2, теорема 2.1]{KozhKhad:bib:Lions}

существует $u\in \stackrel {\circ}{W}\!\!_{\bf B}^{1}(\Omega)$ такая, что $\mathbf{A}(u)=\mathbf{0}$.

Таким образом, для любого $v\in\stackrel {\circ}{W}\!_{\bf B}^{1}(\Omega)$ справедливо тождество~\eqref{eq:Kozh:intt}.

$\square$





\smallskip







\Section{Ограниченность решения}

Доказательство теоремы \hyperlink{kozh:th2}{2} основано на применении следующей леммы \cite{KozhKhad:bib:Korolev2}.



\smallskip



\hypertarget{kozh:lemma4}{}

\begin{lemma}[4]Пусть $N$-функция $B^*(z)$ удовлетворяет $\Delta_2$-условию,  $\mu(t)$ "--- неотрицательная невозрастающая функция на $[a,\infty)$ такая{\rm ,} что

$$\mu(h)\leq cg_0(\mu(k))/B^*(h-k), \quad h>k>a,$$ где $g_0(t),$ $t\geq 0 $ "--- функция{\rm ,} удовлетворяющая условиям

$$g_0(zt)\leq g_0(z)g_0(t),\quad t/g_0(t) \to \infty \;\mbox {при }\; t \to  0.$$

Тогда существует число $\varrho>0$ такое{\rm ,} что $\mu(a+\varrho)=0.$

\end{lemma}



\smallskip



{\it Д\,о\,к\,а\,з\,а\,т\,е\,л\,ь\,с\,т\,в\,о \ т\,е\,о\,р\,е\,м\,ы \rm \hyperlink{kozh:th2}{2}.}

Для $k>0$ определим функцию

 $$u^{(k)}=\mathop{\rm sign}u\max(|u|-k, 0).$$

Из принадлежности  $u\in \stackrel {\circ}{W}\!\!_{\bf B}^{1}(\Omega)$ следует принадлежность  $u^{(k)}\in \stackrel {\circ}{W}\!\!_{\bf B}^{1}(\Omega)$, кроме того, $u_{x_{\alpha}}=u^{(k)}_{x_{\alpha}}$ на множестве $E_k=\{{\bf x}\in \Omega\,\mid\,|u({\bf x})|\geq k\}$.

Пользуясь   \eqref{eq:Kozh:the_3}, выводим

$$

\mathop{\rm mes}  E_k {B_0(k)}\leq  \int_{E_k}B_0(u)d{\bf x}\leq  \int_{\Omega}B_0(u)d{\bf x}\leq C_1,

$$

откуда следует

$$

\mathop{\rm mes} E_k\leq \frac{C_1}{B_0(k)}.$$



Из определения множества $E_k$ имеем $u^{(k)}({\bf x})=0$ в  $\mathbb{R}_n\setminus E_k$. В \eqref{eq:Kozh:intt} положим $v=u^{(k)}$, тогда получим

$$

({\bf A}(u),u^{(k)})_{E_k}=0.

$$

Из условия \eqref{eq:Kozh:us4} следует неравенство

$$(a_0({\bf x},u,\nabla u)-a_0({\bf x},u^{(k)},\nabla u^{(k)}))(u-u^{(k)})\geq 0\quad \text{ для~~ ${\bf x}\in E_k,$}$$  из которого несложно установить

$$a_0({\bf x},u,\nabla u)u^{(k)}\geq a_0({\bf x},u^{(k)},\nabla u^{(k)})u^{(k)}.$$

Кроме того, из \eqref{eq:Kozh:us5} вытекает

$$\sum \limits_{\alpha=1}^n a_{\alpha}({\bf x},u,\nabla u^{(k)})u^{(k)}_{x_{\alpha}}\geq\sum \limits_{\alpha=1}^n a_{\alpha}({\bf x},u^{(k)},\nabla u^{(k)})u^{(k)}_{x_{\alpha}}.$$

Таким образом, справедливо неравенство

$$

({\bf A}(u^{(k)}),u^{(k)})_{E_k}\leq 0.

$$



Далее, применяя \eqref{eq:Kozh:us1}, выводим соотношение

$$

\overline{a}\|\mathbf{B}(u^{(k)})\|_{1,E_k}

\leq\|\psi\|_{1,E_k},

$$

из которого, согласно \eqref{eq:Kozh:the_2_1}, получаем

\begin{multline}

\label{eq:Kozh:the_2(2)}

\int _{E_k}\sum_{\alpha=1}^nB_\alpha(u^{(k)}_{x_\alpha})d{\bf x}\leq\|\mathbf{B}(u^{(k)})\|_{1,E_k}\leq

\\

 \leq

 C_2\|\psi\|_{l/(l-1)}(\mathop{\rm mes} E_k)^{1/l}=C_3(\mathop{\rm mes} E_k)^{1/l}.

\end{multline}



Имеет место следующая цепочка неравенств: $\forall h>k\; E_h\subset E_k$ и

$$B^*(h-k)\mathop{\rm mes} E_h\leq \int_ {E_h}B^*(\mid u\mid-k)d{\bf x}=\int_ {E_h}B^*(u^{(k)})d{\bf x}\leq\int_ {E_k}B^*(u^{(k)})d{\bf x}.$$

Тогда из \eqref{eq:Kozh:the_2(2)}, применяя \eqref{eq:Kozh:K3},  выводим

$$

B^*(h-k)\mathop{\rm mes} E_h\leq\chi(C_4 (\mathop{\rm mes} E_k)^{1/l})\leq C_5\chi( (\mathop{\rm mes} E_k)^{1/l}).$$

Утверждение теоремы вытекает из  леммы \hyperlink{kozh:lemma4}{4}, применённой к функциям  $\mu(t)=\mathop{\rm mes} E_t$ и $g_0(t)=\chi(t^{1/l})=t^{1/l}\Lambda(t^{1/(ln)})$.

$\square$



\smallskip









\Section{Убывание решения}

Далее будем предполагать, что область $\Omega$  расположена вдоль выделенной оси

$Ox_s$, $s\in\{1, 2,\ldots,n\}$.

Для функций $v({\bf x})\in\stackrel {\circ}

{H}{}_{B}^1(\Omega)$ имеет место аналог  неравенства Фридрихса:

\begin{equation}\label{eq:Kozh:frid}

\int_{\Omega^r}B_s\Bigl(\frac{v({\bf x})}{r}\Bigr)d{\bf x}\leq

\int_{\Omega^r}B_s(v_{x_s}({\bf x}))d{\bf x},\quad r>0,

\end{equation}

(см. \cite [неравенство (57)]{KozhKhad:bib:Andrianova}).



\smallskip



{\it Д\,о\,к\,а\,з\,а\,т\,е\,л\,ь\,с\,т\,в\,о \ т\,е\,о\,р\,е\,м\,ы \rm \hyperlink{kozh:th3}{3}.}

Пусть $\theta_+(x)$, $x>0$ "--- абсолютно непрерывная функция,

равная единице при $x\ge r$, нулю при $x\le R_0$, линейная

при $x\in[R_0, 2R_0]$ и удовлетворяющая уравнению

\begin{equation}\label{eq:Kozh:the_4(0)}

\theta'_+(x)=\delta \overline{\nu}_s^*(x)\theta_+(x),\quad

x\in(2R_0,r),

\end{equation}

 (постоянную $\delta$ определим позднее).   Решая это уравнение, находим

 $$

\theta_+(x)=\exp\left(-\delta \int_{x}^r

\overline{\nu}_s^*(\rho)d\rho\right),\quad x\in(2R_0,r),

$$

 тогда имеем

\begin{equation}

\theta'_+(x)=\frac{\theta_+(2R_0)}{R_0}=\frac{1}{R_0}\exp\left(-\delta

\int_{2R_0}^r \overline{\nu}_s^*(\rho)d\rho\right),\quad

x\in(R_0,2R_0).\label{eq:Kozh:the_4(1)}

\end{equation}



Подставив в \eqref{eq:Kozh:intt} $v =u\theta_+(x_s)$, получим

\begin{multline*}

  \int_{\Omega}\Biggl\{\biggl(\sum\limits_{\alpha=1}^{n}a_{\alpha}({\bf

x},u,\nabla u)u_{x_{\alpha}}+a_0({\bf

x}, u,\nabla u)u\biggr)\theta_+(x_s)+ \\

+ a_{s}({\bf x,}u,\nabla u)u\theta'_+(x_s)\Biggr\}d{\bf x}=0.

\end{multline*}



Используя \eqref{eq:Kozh:us1}, \eqref{eq:Kozh:the_4(0)}, \eqref {eq:Kozh:the_4(1)}, \eqref{eq:Kozh:supp} и учитывая,

что носители функций $\theta$ и $\psi$, $\psi_1$  не пересекаются, получим

\begin{multline}\label{eq:Kozh:the_4(3)}

  \overline{a}\int_{\Omega}\theta_+(x_s)\mathbf{B}(u)d{\bf x}

    \leq \int_{\Omega_{2R_0}^r} |u| |a_{s}({\bf x},u,\nabla u)| \delta

  \overline{\nu}_s^*(x_s)\theta_+(x_s)d{\bf x}+

\\

+\int_{\Omega_{R_0}^{2R_0}}|u| \, |a_{s}({\bf x},u,\nabla

  u)|\frac{\theta_+(2R_0)}{R_0}d{\bf x}=I_1+I_2.

\end{multline}

 Далее, пользуясь неравенством $\overline{\nu}_s^*(x_s)\leq z^*$ и

ограниченностью функции  $u({\bf x})$ (см. \eqref{eq:Kozh:the_2_2}), при помощи \eqref{eq:Kozh:us2}, \eqref{eq:Kozh:BB} в предположении  $\frac{\delta}{\varepsilon} \mathcal{M}_2\leq 1$ оценим

первый интеграл:

\begin{multline*}

  I_1 \leq \int_{\Omega_{2R_0}^r}

  \theta_+(x_s)\left(\overline{B}_s(\varepsilon a_{s}({\bf x},u,\nabla

  u))+B_s\Bigl(u\overline{\nu}_s^*(x_s)\frac{\delta}{\varepsilon}\Bigr)\right)d{\bf x}\leq

\\

  \leq \int_{\Omega_{2R_0}^r} \theta_+(x_s)\biggl(\varepsilon \widehat{a}_1\mathbf{B}(u)+\sum\limits_{i=1,i\neq s}^n \beta_i B_i\Bigl(u{\nu}_i(x_s)\frac{\delta}{\varepsilon}\Bigr)\biggr)d{\bf x}.

\end{multline*}

Выберем теперь $\varepsilon={\overline{a}}/({4\widehat{a}_1})$, а также $\delta$ так,

чтобы

$$\frac{\delta}{\varepsilon}\leq \max\{1,\mathcal{M}^{-1}_2\}, \quad \frac{\delta}{\varepsilon}\max \limits_{i=\overline{1,n},i\neq s}\beta_i\leq \frac{\overline{a}}{4}.$$ Тогда, воспользовавшись

определением \eqref{eq:Kozh:nu} функций ${\nu}_i$, получим

\begin{multline*}

  I_1\leq \frac{\overline{a}}{4}\int_{\Omega_{2R_0}^r}\theta_+ (x_s)\biggl(\mathbf{B}(u)+ \sum\limits_{i=1,i\neq s}^n B_i(u\nu_i)\biggr)d{\bf x}\leq

\\

 \leq \frac{\overline{a}}{4}\int_{\Omega_{2R_0}^r}\theta_+(x_s) \biggl(\mathbf{B}(u)+

 \sum\limits_{i=1,i\neq s}^n B_i(u_{x_i})\biggr)d{\bf x}\leq

\frac{\overline{a}}{2}

  \int_{\Omega_{2R_0}^r} \theta_+(x_s)\mathbf{B}(u) d{\bf x}.

\end{multline*}



  Теперь, используя условие \eqref{eq:Kozh:us2} и неравенства  \eqref{eq:Kozh:frid}, \eqref{eq:Kozh:delta2}, выводим

\begin{multline*}

  I_2\leq{\theta_+(2R_0)}

  \int_{\Omega_{R_0}^{2R_0}}\left(B_s\Bigl(\frac{u}{R_0}\Bigr)+\overline{B}_s(a_{s}({\bf x},u,\nabla

  u))\right)d{\bf x}\leq

\\

  \leq{\theta_+(2R_0)}

  \int_{\Omega^{2R_0}}\bigl(B_s(2u_{x_s})+\widehat{a}_1\mathbf{B}(u)\bigr)d{\bf x}\leq{\theta_+(2R_0)}C_1 \|\mathbf{B}(u)\|_{1,\Omega^{2R_0}}.

\end{multline*}

Подставляя в \eqref{eq:Kozh:the_4(3)} оценки для $I_1, I_2$ и применяя \eqref{eq:Kozh:the_3}, выводим \eqref{eq:Kozh:the_4}.

$\square$



\thanks{\textbf{Благодарности.} Работа выполнена при финансовой поддержке Российского фонда фундаментальных исследований (проект \No~13--01--00081-a).}





\thanks{{{\bf ORCID}







{\bf Лариса Михайловна Кожевникова:} \url{http://orcid.org/0000-0002-6458-5998}



{\bf Анна Александровна Хаджи:} \url{http://orcid.org/0000-0003-0979-5893}



}}









\RusRef



\begin{thebibliography}{20}



\RBibitem{kozh:add1-ru}

\by Кожевникова~Л.~М., Хаджи~А.~А.

\paper  О решениях эллиптических уравнений с~нестепенными нелинейностями в неограниченных областях

\pages 199--200

\inbook Четвертая международная конференция «Математическая физика и ее приложения»

\bookinfo материалы конф.

\eds чл.-корр.~РАН И.~В.~Волович; д.ф.-м.н., проф. В.~П.~Радченко

\publaddr Самара

\publ СамГТУ

\yr 2014





\RBibitem{KozhKhad:bib:Vichik}

\by Вишик~М.~И.

\paper О~разрешимости первой краевой задачи для квазилинейных уравнений с быстрорастущими коэффициентами в классах Орлича

\jour ДАН СССР

\yr 1963

\vol 151

\issue 4

\pages 758--761

\zmath{https://zbmath.org/?q=an:03254321}

%Solvability of the first boundary value problem for quasilinear equations with rapidly increasing coefficients in Orlicz classes. (English. Russian original) Zbl 0158.12403

%Sov. Math., Dokl. 4, 1060-1064 (1963); translation from Dokl. Akad. Nauk SSSR 151, 758-761 (1963).





\RBibitem{KozhKhad:bib:Dubinskii}

\by Дубинский~Ю.~А.

\paper Слабая сходимость в~нелинейных эллиптических и параболических уравнениях

\jour Матем. сб.

\yr 1965

\vol 67(109)

\issue 4

\pages 609--642

\mathnet{http://mi.mathnet.ru/msb4394}

\mathscinet{http://www.ams.org/mathscinet-getitem?mr=190546}

\zmath{https://zbmath.org/?q=an:0145.35202}



\Bibitem{KozhKhad:bib:Donaldson}

\by  Donaldson~T.

\paper Nonlinear elliptic boundary value problems in Orlicz--Sobolev spaces

\jour J. Diff. Eq.

\yr 1971

\vol 10

\issue 3

\pages 507--528

\crossref{10.1016/0022-0396(71)90009-x}



\RBibitem{KozhKhad:bib:Klimov2}

\by Климов~В.~С.

\inbook  Качественные и приближенные методы исследования операторных уравнений

%Qualitative methods for investigating operator equation

\paper Краевые задачи в пространствах Орлича--Соболева

%Boundary value problems in Orlicz -Sobolev spaces

\publaddr Ярославль

\publ Ярославский государственный университет

\yr 1976

\pages 75--93







\Bibitem{KozhKhad:bib:Giorgi}

\lang In Italian

\by  De Giorgi~E.

\paper Sulla differenziabilit\`a e l'analiticit\`a delle estremali degli integrali multipli regolari

\jour Mem. Acad. Sci. Torino, Serie III

\yr 1957

\vol 3

\pages 25--43

\zmath{https://zbmath.org/?q=an:0084.31901}

\mathscinet{http://www.ams.org/mathscinet-getitem?mr=0227827}

\transl

\by  De~Giorgi~E.

\paper On the differentiability and the analiticity of extremals of regular multiple integrals

\pages 149--166

\yr 2006

\inbook Selected papers

\publaddr Berlin, New York

\eds Luigi Ambrosio, Gianni Dal Maso, Marco Forti, Mario Miranda, and Sergio Spagnolo

\publ Springer-Verlag

\zmath{https://zbmath.org/?q=an:1096.01015}

\mathscinet{www.ams.org/mathscinet-getitem?mr=2229237}





\Bibitem{KozhKhad:bib:Moser}

\by  Moser~J.

\paper A new proof of de Giorgi's theorem concerning the regularity problem for elliptic differential equations

\jour Comm. Pure Appl. Math.

\yr 1960

\vol 13

\issue 3

\pages 457--468

\crossref{10.1002/cpa.3160130308}





\RBibitem{KozhKhad:bib:Kruzhkov}

\by Кружков~С.~Н.

\paper Априорные оценки и некоторые свойства решений эллиптических и параболических уравнений

\jour Матем. сб.

\yr 1964

\vol 65(107)

\issue 4

\pages 522--570

\mathnet{http://mi.mathnet.ru/msb4494}

\mathscinet{http://www.ams.org/mathscinet-getitem?mr=171086}

\zmath{https://zbmath.org/?q=an:0148.35703}







\Bibitem{KozhKhad:bib:Serrin}

\by  Serrin~J.

\paper Local behavior of solutions of quasi-linear equations

\jour Acta Math.

\yr 1964

\vol 111

\issue 1

\pages 247--302

\crossref{10.1007/BF02391014}





\RBibitem{KozhKhad:bib:Landis}

\by Ландис~Е.~М.

\paper Новое доказательство теоремы E.\,Де~Джорджи

\serial Тр. ММО

\yr 1967

\vol 16

\pages 319--328

\publ Издательство Московского университета

\publaddr М.

\mathnet{http://mi.mathnet.ru/mmo188}

\mathscinet{http://www.ams.org/mathscinet-getitem?mr=224975}

\zmath{https://zbmath.org/?q=an:0182.43403}





\RBibitem{KozhKhad:bib:Kolodii}

\by Колодий~И.~М.

\paper Об ограниченности обобщенных решений эллиптических дифференциальных уравнений

\jour Вестн. Моск. унив., Сер.~1.

\yr 1970

\issue 5

\pages 45--52

\zmath{https://zbmath.org/?q=an:0207.10701}

%Kolodij, I.M.

%The boundedness of generalized solutions of elliptic differential equations. (English) Zbl 0245.35029

%Mosc. Univ. Math. Bull. 25(1970), No.5-6, 31-37 (1972).

%\zmath{https://zbmath.org/?q=an:0245.35029}





\RBibitem{KozhKhad:bib:KozhKhad2}

\by Кожевникова~Л.~М., Хаджи~А.~А.

\paper Ограниченность решений анизотропных эллиптических уравнений второго порядка в~неограниченных областях

\jour Уфимск. матем. журн.

\yr 2014

\vol 6

\issue 2

\pages 67--77

\mathnet{http://mi.mathnet.ru/ufa244}

%\transl

%\jour Ufa Math. Journal

%\yr 2014

%\vol 6

%\issue 2

%\pages 66--76

%\crossref{10.13108/2014-6-2-66}

%\scopus{http://www.scopus.com/record/display.url?origin=inward&eid=2-s2.0-84928198601}





\RBibitem{KozhKhad:bib:Korolev3}

\by Королев~А.~Г.

\paper Об ограниченности обобщенных решений

эллиптических дифференциальных уравнений

с нестепенными нелинеиностями

\jour Матем. сб.

\yr 1989

\vol 180

\issue 1

\pages 78--100

\mathnet{http://mi.mathnet.ru/msb1598}

\mathscinet{http://www.ams.org/mathscinet-getitem?mr=988847}

\zmath{https://zbmath.org/?q=an:0706.35030}

%\transl

%\by A.~G.~Korolev

%\paper On boundedness of generalized solutions of elliptic differential equations with nonpower nonlinearities

%\jour Math. USSR-Sb.

%\yr 1990

%\vol 66

%\issue 1

%\pages 83--106

%\crossref{10.1070/SM1990v066n01ABEH001166}

%\isi{http://gateway.isiknowledge.com/gateway/Gateway.cgi?GWVersion=2&SrcApp=PARTNER_APP&SrcAuth=LinksAMR&DestLinkType=FullRecord&DestApp=ALL_WOS&KeyUT=A1990DK06800004}







\RBibitem{KozhKhad:bib:LadizhYralc}

\by Ладыженская~О.~А., Уральцева~Н.~Н.

\book Линейные и квазилинейные уравнения эллиптического типа

\publaddr М.

\publ Наука

\yr 1973

\totalpages 578

\zmath{https://zbmath.org/?q=an:03424452}

%Ladyzhenskaya, O.A.; Ural’tseva, N.N.

%Linear and quasilinear elliptic equations. (English) Zbl 0164.13002

%Mathematics in Science and Engineering. 46. New York-London: Academic Press. XVIII, 495 p. (1968).

%\zmath{https://zbmath.org/?q=an:03262653}





\RBibitem{KozhKhad:bib:Klimov3}

\by Климов~В.~С.

\paper Теоремы вложения для пространств Орлича и~их приложения  к~краевым задачам

\jour Сиб. матем. журн.

\yr 1972

\vol 13

\pages 334--348

\zmath{https://zbmath.org/?q=an:03379592}





\RBibitem{KozhKhad:bib:Korolev2}

\by Королев~А.~Г.

\paper Об~ограниченности обобщенных решений эллиптических

уравнений с~нестепенными нелинейностями

\jour Матем. заметки

\yr 1987

\vol 42

\issue 2

\pages 244--255

\mathnet{http://mi.mathnet.ru/mz4979}

\mathscinet{http://www.ams.org/mathscinet-getitem?mr=915112}

\zmath{https://zbmath.org/?q=an:0681.35022}

%\transl

%\jour Math. Notes

%\yr 1987

%\vol 42

%\issue 2

%\pages 639--645

%\crossref{http://dx.doi.org/10.1007/BF01240452}

%\isi{http://gateway.isiknowledge.com/gateway/Gateway.cgi?GWVersion=2&SrcApp=PARTNER_APP&SrcAuth=LinksAMR&DestLinkType=FullRecord&DestApp=ALL_WOS&KeyUT=A1987N113800026}







\RBibitem{KozhKhad:bib:OlIosif}

\by Олейник~О.~А., Иосифьян~Г.~А.

\paper О~поведении на бесконечности решений эллиптических уравнений второго порядка в~областях с~некомпактной границей

\jour Матем. сб.

\yr 1980

\vol 112(154)

\issue 4(8)

\pages 588--610

\mathnet{http://mi.mathnet.ru/msb2738}

\mathscinet{http://www.ams.org/mathscinet-getitem?mr=587039}

\zmath{https://zbmath.org/?q=an:0469.35045|0445.35038}

%\transl

%\by O.~A.~Oleinik, G.~A.~Iosif'yan

%\paper On the behavior at infinity of solutions of second order elliptic equations in domains with noncompact boundary

%\jour Math. USSR-Sb.

%\yr 1981

%\vol 40

%\issue 4

%\pages 527--548

%\crossref{10.1070/SM1981v040n04ABEH001849}

%\isi{http://gateway.isiknowledge.com/gateway/Gateway.cgi?GWVersion=2&SrcApp=PARTNER_APP&SrcAuth=LinksAMR&DestLinkType=FullRecord&DestApp=ALL_WOS&KeyUT=A1981MW11700004}







\RBibitem{KozhKhad:bib:KondKopLek}

\by Кондратьев~В.~А.,  Копачек~И.,  Леквеишвили~Д.~М., Олейник~О.~А.

\paper Неулучшаемые оценки в~пространствах Гельдера и точный принцип Сен-Венана для решений бигармонического уравнения

\inbook Современные проблемы математики. Дифференциальные уравнения, математический анализ и их приложения

\bookinfo Сборник статей. Посвящается академику Льву Семеновичу Понтрягину к~его семидесятипятилетию

\serial Тр. МИАН СССР

\yr 1984

\vol 166

\pages 91--106

\mathnet{http://mi.mathnet.ru/tm2254}

\mathscinet{http://www.ams.org/mathscinet-getitem?mr=752171}

\zmath{https://zbmath.org/?q=an:0566.35045|0584.35038}

%\transl

%\jour Proc. Steklov Inst. Math.

%\yr 1986

%\vol 166

%\pages 97--116



\RBibitem{KozhKhad:bib:Kozhevnikova}

\by Кожевникова~Л.~М.

\paper Поведение на бесконечности решений псевдодифференциальных

эллиптических уравнений в~неограниченных областях

\jour Матем. сб.

\yr 2008

\vol 199

\issue 8

\pages 61--94

\mathnet{http://mi.mathnet.ru/msb4235}

\crossref{10.4213/sm4235}

\mathscinet{http://www.ams.org/mathscinet-getitem?mr=2452267}

%\transl

%\by L.~M.~Kozhevnikova

%\paper Behaviour at infinity of solutions of pseudodifferential

%elliptic equations in unbounded domains

%\jour Sb. Math.

%\yr 2008

%\vol 199

%\issue 8

%\pages 1169--1200

%\crossref{10.1070/SM2008v199n08ABEH003958}

%\isi{http://gateway.isiknowledge.com/gateway/Gateway.cgi?GWVersion=2&SrcApp=PARTNER_APP&SrcAuth=LinksAMR&DestLinkType=FullRecord&DestApp=ALL_WOS&KeyUT=000260697900011}

%\scopus{http://www.scopus.com/record/display.url?origin=inward&eid=2-s2.0-57049087573}







\RBibitem{KozhKhad:bib:GilMuk}

\by Гилимшина~В.~Ф., Мукминов~Ф.~Х.

\paper Об убывании решения неравномерно эллиптического уравнения

\jour Изв. РАН. Сер. матем.

\yr 2011

\vol 75

\issue 1

\pages 53--70

\mathnet{http://mi.mathnet.ru/izv3292}

\crossref{10.4213/im3292}

\mathscinet{http://www.ams.org/mathscinet-getitem?mr=2815995}

\zmath{https://zbmath.org/?q=an:1210.35131}

\adsnasa{http://adsabs.harvard.edu/cgi-bin/bib_query?2011IzMat..75...53G}

%\transl

%\jour Izv. Math.

%\yr 2011

%\vol 75

%\issue 1

%\pages 53--71

%\crossref{10.1070/IM2011v075n01ABEH002527}

%\isi{http://gateway.isiknowledge.com/gateway/Gateway.cgi?GWVersion=2&SrcApp=PARTNER_APP&SrcAuth=LinksAMR&DestLinkType=FullRecord&DestApp=ALL_WOS&KeyUT=000287579900003}

%\scopus{http://www.scopus.com/record/display.url?origin=inward&eid=2-s2.0-80053475855}



\RBibitem{KozhKhad:bib:KozhKarimov}

\by Кожевникова~Л.~М., Каримов~Р.~Х.

\paper Поведение на бесконечности решений квазилинейных эллиптических уравнений второго порядка в~неограниченных областях

\jour Уфимск. матем. журн.

\yr 2010

\vol 2

\issue 2

\pages 53--66

\mathnet{http://mi.mathnet.ru/ufa51}

\zmath{https://zbmath.org/?q=an:1240.35043}



\RBibitem{KozhKhad:bib:KozhKhad1}

\by Кожевникова~Л.~М., Хаджи~А.~А.

\paper Решения анизотопных эллиптических уравнений в~неограниченных областях

\jour Вестн. Сам. гос. техн. ун-та. Сер. Физ.-мат. науки

\yr 2013

\vol 1(30)

\pages 90--96

\mathnet{http://mi.mathnet.ru/vsgtu1163}

\crossref{10.14498/vsgtu1163}





\RBibitem{KozhKhad:bib:Rut}

\by Рутицкий~Я.~Б., Красносельский~М.~А.

\book Выпуклые функции и пространства Орлича

\publaddr М.

\publ  Физматлит

\yr 1958

\totalpages 587

\zmath{https://zbmath.org/?q=an:03137657}

%Krasnosel’skij, M.A.; Rutitskij, Ya.B.

%Convex functions and Orlicz spaces. (English. Russian original) Zbl 0095.09103

%Groningen-The Netherlands: P. Noordhoff Ltd. IX, 249 p. (1961).

%\zmath{https://zbmath.org/?q=an:03155055}





\RBibitem{KozhKhad:bib:Korolev1}

\paper Теоремы вложения анизотропных пространств Соболева--Орлича

\by Королев~А.~Г.

\jour Вестн. Моск. унив., Сер.~1.

\yr 1983

\issue 1

\pages 32--37

\zmath{https://zbmath.org/?q=an:03820448}

%Korolev, A.G.

%Embedding theorems for anisotropic Sobolev-Orlicz spaces. (English. Russian original) Zbl 0518.46023

%Mosc. Univ. Math. Bull. 38, No.1, 37-42 (1983); translation from Vestn. Mosk. Univ., Ser. I 1983, No.1, 32-37 (1983).





\RBibitem{KozhKhad:bib:Lions}

\by Лионс~Ж.~Л.

\book Некоторые методы решения нелинейных краевых задач

\publaddr М.

\publ  Мир

\yr 1972

\totalpages 596

%Lions, J.L.

%Some methods in the mathematical analysis of systems and their control. (English) Zbl 0542.93034

%Beijing, China: Science Press; New York: Gordon and Breach, Science Publishers, Inc. XXIII, 542 p. $ 93.50 (1981).



\RBibitem{KozhKhad:bib:Andrianova}

\by Андриянова~Э.~Р.

\paper Оценки скорости убывания решения параболического уравнения с~нестепенными нелинейностями

\jour Уфимск. матем. журн.

\yr 2014

\vol 6

\issue 2

\pages 3--25

\mathnet{http://mi.mathnet.ru/ufa239}

%\transl

%\jour Ufa Math. Journal

%\yr 2014

%\vol 6

%\issue 2

%\pages 3--24

%\crossref{10.13108/2014-6-2-3}

%\scopus{http://www.scopus.com/record/display.url?origin=inward&eid=2-s2.0-84928186140}



\end{thebibliography}



\RuDate





\newpage





\makeentitle





\thanks{\textbf{Acknowledgments.}

This work has been  supported by the Russian Foundation for Basic Research

 (project no.~13--01--00081-a).}





\thanks{{{\bf ORCID}







{\bf Larisa M. Kozhevnikova:} \url{http://orcid.org/0000-0002-6458-5998}



{\bf Anna A. Khadzhi:} \url{http://orcid.org/0000-0003-0979-5893}



}}











\EngRef

\English



\begin{thebibliography}{20}



\Bibitem{kozh:add1-eng}

\lang In Russian

\by Kozhevnikova~L.~M., Khadzhi~A.~A.

\paper On solutions of elliptic equations with nonpower nonlinearities in unbounded domains

\pages 199-200

\inbook  The 4nd International Conference ``Mathematical Physics and its Applications''

\bookinfo Book of Abstracts and Conference Materials

\eds I.~V.~Volovich; V.~P.~Radchenko

\publaddr Samara

\publ Samara State Technical Univ.

\yr 2014







\Bibitem{KozhKhad:bib:Vichik:eng}

\zmath{https://zbmath.org/?q=an:03254321}

\lang In Russian

\by Vishik~M.~I.

\paper Solvability of the first boundary value problem for quasilinear equations with rapidly increasing coefficients in Orlicz classes

\jour Sov. Math., Dokl.

\vol 4

\pages  1060--1064

\yr 1963







\Bibitem{KozhKhad:bib:Dubinskii:eng}

\lang In Russian

\by Dubinskii~Yu.~A.

\paper Weak convergence for nonlinear elliptic and parabolic equations

\jour Mat. Sb. (N.S.)

\yr 1965

\vol 67(109)

\issue 4

\pages 609--642



\Bibitem{KozhKhad:bib:Donaldson:eng}

\by  Donaldson~T.

\paper Nonlinear elliptic boundary value problems in Orlicz--Sobolev spaces

\jour J. Diff. Eq.

\yr 1971

\vol 10

\issue 3

\pages 507--528

\crossref{10.1016/0022-0396(71)90009-x}



\Bibitem{KozhKhad:bib:Klimov2:eng}

\lang In Russian

\by Klimov~V.~S.

\inbook  Kachestvennye i priblizhennye metody issledovaniia operatornykh uravnenii \rm [Qualitative methods for investigating operator equation]

\paper Boundary value problems in Orlicz--Sobolev spaces

\publaddr Yaroslavl

\publ Yaroslavl State Univ.

\yr 1976

\pages 75--93





\Bibitem{KozhKhad:bib:Giorgi:eng}

\lang In Italian

\by  De Giorgi~E.

\paper Sulla differenziabilit\`a e l'analiticit\`a delle estremali degli integrali multipli regolari

\jour Mem. Acad. Sci. Torino, Serie III

\yr 1957

\vol 3

\pages 25--43

\zmath{https://zbmath.org/?q=an:0084.31901}

\mathscinet{http://www.ams.org/mathscinet-getitem?mr=0227827}

\transl

\by  De~Giorgi~E.

\paper On the differentiability and the analiticity of extremals of regular multiple integrals

\pages 149--166

\yr 2006

\inbook Selected papers

\publaddr Berlin, New York

\eds Luigi Ambrosio, Gianni Dal Maso, Marco Forti, Mario Miranda, and Sergio Spagnolo

\publ Springer-Verlag

\zmath{https://zbmath.org/?q=an:1096.01015}

\mathscinet{www.ams.org/mathscinet-getitem?mr=2229237}





\Bibitem{KozhKhad:bib:Moser:eng}

\by  Moser~J.

\paper A new proof of de Giorgi's theorem concerning the regularity problem for elliptic differential equations

\jour Comm. Pure Appl. Math.

\yr 1960

\vol 13

\issue 3

\pages 457--468

\crossref{10.1002/cpa.3160130308}





\Bibitem{KozhKhad:bib:Kruzhkov:eng}

\lang In Russian

\by Kruzhkov~S.~N.

\paper A priori bounds and some properties of solutions of elliptic and parabolic equations

\jour Mat. Sb. (N.S.)

\yr 1964

\vol 65(107)

\issue 4

\pages 522--570









\Bibitem{KozhKhad:bib:Serrin:eng}

\by  Serrin~J.

\paper Local behavior of solutions of quasi-linear equations

\jour Acta Math.

\yr 1964

\vol 111

\issue 1

\pages 247--302

\crossref{10.1007/BF02391014}





\Bibitem{KozhKhad:bib:Landis:eng}

\lang In Russian

\by Landis~E.~M.

\paper A new proof of E.\,De~Giorgi's theorem

\serial Tr. Mosk. Mat. Obs.

\yr 1967

\vol 16

\pages 319--328

\publ MSU

\publaddr Moscow





\Bibitem{KozhKhad:bib:Kolodii:eng}

\by Kolodij~I.~M.

\paper The boundedness of generalized solutions of elliptic differential equations

\jour Mosc. Univ. Math. Bull.

\vol 25

\yr 1970

\issue 5-6

\pages 31--37

\yr 1972

\zmath{https://zbmath.org/?q=an:0245.35029}





\Bibitem{KozhKhad:bib:KozhKhad2:eng}

\paper Boundedness of solutions to anisotropic second order elliptic equations in unbounded domains

\by Kozhevnikova~L.~M., Khadzhi~A.~A.

\jour Ufa Math. Journal

\yr 2014

\vol 6

\issue 2

\pages 66--76

\crossref{10.13108/2014-6-2-66}

\scopus{http://www.scopus.com/record/display.url?origin=inward&eid=2-s2.0-84928198601}





\Bibitem{KozhKhad:bib:Korolev3:eng}

\by Korolev~A.~G.

\paper On boundedness of generalized solutions of elliptic differential equations with nonpower nonlinearities

\jour Math. USSR-Sb.

\yr 1990

\vol 66

\issue 1

\pages 83--106

\crossref{10.1070/SM1990v066n01ABEH001166}

\isi{http://gateway.isiknowledge.com/gateway/Gateway.cgi?GWVersion=2&SrcApp=PARTNER_APP&SrcAuth=LinksAMR&DestLinkType=FullRecord&DestApp=ALL_WOS&KeyUT=A1990DK06800004}







\Bibitem{KozhKhad:bib:LadizhYralc:eng}

\by Ladyzhenskaya~O.~A., Ural'tseva~N.~N.

\book Linear and quasilinear elliptic equations

\serial Mathematics in Science and Engineering

\vol 46

\publaddr New York, London

\publ Academic Press

\totalpages xviii+495

\yr 1968

\zmath{https://zbmath.org/?q=an:03262653}





\Bibitem{KozhKhad:bib:Klimov3:eng}

\by Klimov~V.~S.

\paper Imbedding theorems for Orlicz spaces and their applications to boundary value problems

\jour Sib. Math.~J.

\yr 1972

\vol 13

\pages 231-240

\zmath{https://zbmath.org/?q=an:0246.46022}

\crossref{10.1007/BF00971611}







\Bibitem{KozhKhad:bib:Korolev2:eng}

\paper Boundedness of the generalized solutions of elliptic equations with nonpower nonlinearities

\by Korolev~A.~G.

\jour Math. Notes

\yr 1987

\vol 42

\issue 2

\pages 639--645

\crossref{10.1007/BF01240452}

\isi{http://gateway.isiknowledge.com/gateway/Gateway.cgi?GWVersion=2&SrcApp=PARTNER_APP&SrcAuth=LinksAMR&DestLinkType=FullRecord&DestApp=ALL_WOS&KeyUT=A1987N113800026}







\Bibitem{KozhKhad:bib:OlIosif:eng}

\by Oleinik~O.~A., Iosif'yan~G.~A.

\paper On the behavior at infinity of solutions of second order elliptic equations in domains with noncompact boundary

\jour Math. USSR-Sb.

\yr 1981

\vol 40

\issue 4

\pages 527--548

\crossref{10.1070/SM1981v040n04ABEH001849}

\isi{http://gateway.isiknowledge.com/gateway/Gateway.cgi?GWVersion=2&SrcApp=PARTNER_APP&SrcAuth=LinksAMR&DestLinkType=FullRecord&DestApp=ALL_WOS&KeyUT=A1981MW11700004}







\Bibitem{KozhKhad:bib:KondKopLek:eng}

\jour Proc. Steklov Inst. Math.

\yr 1986

\vol 166

\pages 97--116

\by Kondrat'ev~V.~A., Kop\'a\v{c}ek~J., Lekveishvili~D.~M., Ole\u{\i}nik~O.~A.

\paper Sharp estimates in H\"older spaces and precise Saint-Venant principle for solutions of the biharmonic equation



\Bibitem{KozhKhad:bib:Kozhevnikova:eng}

\by Kozhevnikova~L.~M.

\paper Behaviour at infinity of solutions of pseudodifferential

elliptic equations in unbounded domains

\jour Sb. Math.

\yr 2008

\vol 199

\issue 8

\pages 1169--1200

\crossref{10.1070/SM2008v199n08ABEH003958}

\isi{http://gateway.isiknowledge.com/gateway/Gateway.cgi?GWVersion=2&SrcApp=PARTNER_APP&SrcAuth=LinksAMR&DestLinkType=FullRecord&DestApp=ALL_WOS&KeyUT=000260697900011}

\scopus{http://www.scopus.com/record/display.url?origin=inward&eid=2-s2.0-57049087573}







\Bibitem{KozhKhad:bib:GilMuk:eng}

\paper On the decay of solutions of non-uniformly elliptic equations

\by Gilimshina~V.~F.

\jour Izv. Math.

\yr 2011

\vol 75

\issue 1

\pages 53--71

\crossref{10.1070/IM2011v075n01ABEH002527}

\isi{http://gateway.isiknowledge.com/gateway/Gateway.cgi?GWVersion=2&SrcApp=PARTNER_APP&SrcAuth=LinksAMR&DestLinkType=FullRecord&DestApp=ALL_WOS&KeyUT=000287579900003}

\scopus{http://www.scopus.com/record/display.url?origin=inward&eid=2-s2.0-80053475855}



\Bibitem{KozhKhad:bib:KozhKarimov:eng}

\lang In Russian

\by Karimov~R.~Kh., Kozhevnikova~L.~M.

\paper Behavior on infinity of decision quasilinear elliptical equations in unbounded domain

\jour Ufimsk. Mat. Zh.

\yr 2010

\vol 2

\issue 2

\pages 53--66



\Bibitem{KozhKhad:bib:KozhKhad1:eng}

\lang In Russian

\by Kozhevnikova~L.~M., Khadzhi~A.~A.

\paper Solutions   of anisotropic elliptic equations in~unbounded domains

\jour Vestn. Samar. Gos. Tekhn. Univ. Ser. Fiz.-Mat. Nauki \rm [J. Samara State Tech. Univ., Ser. Phys. \& Math. Sci.]

\yr 2013

\vol 1(30)

\pages 90--96

\crossref{10.14498/vsgtu1163}





\Bibitem{KozhKhad:bib:Rut:eng}

\by Krasnosel'skij~M.~A.; Rutitskij~Ya.~B.

\book Convex functions and Orlicz spaces

\publaddr Groningen

\publ P.~Noordhoff Ltd.

\totalpages ix+249

\yr 1961

\zmath{https://zbmath.org/?q=an:03155055}





\Bibitem{KozhKhad:bib:Korolev1:eng}

\zmath{https://zbmath.org/?q=an:03820448}

\by Korolev~A.~G.

\paper Embedding theorems for anisotropic Sobolev--Orlicz spaces

\jour Mosc. Univ. Math. Bull.

\vol 38

\issue 1

\pages 37--42

\yr 1983





\Bibitem{KozhKhad:bib:Lions:eng}

\by Lions~J.~L.

\book Some methods in the mathematical analysis of systems and their control

\zmath{https://zbmath.org/?q=an:0542.93034}

\publaddr Beijing, China

\publ Science Press

\publaddrr New York

\publl Gordon and Breach, Science Publishers, Inc

\totalpages xxiii+542

\yr 1981



\Bibitem{KozhKhad:bib:Andrianova:eng}

\paper Estimates of decay rate for solution to parabolic equation with non-power nonlinearities

\by Andriyanova~E.~R.

\jour Ufa Math. Journal

\yr 2014

\vol 6

\issue 2

\pages 3--24

\crossref{10.13108/2014-6-2-3}

\scopus{http://www.scopus.com/record/display.url?origin=inward&eid=2-s2.0-84928186140}



\end{thebibliography}



\EngDate



\newpage

\Russian



%\end{document}



